\markdownRendererDocumentBegin
\markdownRendererWarning{The `hybrid` option has been soft-deprecated.}{The `hybrid` option has been soft-deprecated.}{Consider using one of the following better options for mixing TeX and markdown: `contentBlocks`, `rawAttribute`, `texComments`, `texMathDollars`, `texMathSingleBackslash`, and `texMathDoubleBackslash`. For more information, see the user manual at <https://witiko.github.io/markdown/>.}{Consider using one of the following better options for mixing TeX and markdown: `contentBlocks`, `rawAttribute`, `texComments`, `texMathDollars`, `texMathSingleBackslash`, and `texMathDoubleBackslash`. For more information, see the user manual at <https://witiko.github.io/markdown/>.}Berlekamp–Massey 算法是一种用于求数列的最短递推式的算法。给定一个长为 $n$ 的数列，如果它的最短递推式的阶数为 $m$，则 Berlekamp–Massey 算法能够在 $O(nm)$ 时间内求出数列的每个前缀的最短递推式。最坏情况下 $m = O(n)$，因此算法的最坏复杂度为 $O(n^2)$。\markdownRendererInterblockSeparator
{}\markdownRendererSectionBegin
\markdownRendererSectionBegin
\markdownRendererSectionBegin
\markdownRendererHeadingThree{定义}\markdownRendererInterblockSeparator
{}定义一个数列 ${a_0 \dots a_{n - 1} }$ 的递推式为满足下式的序列 ${r_0\dots r_m}$：\markdownRendererParagraphSeparator
{}$\sum_{j = 0} ^ m r_j a_{i - j} = 0, \forall i \ge m$\markdownRendererParagraphSeparator
{}其中 $r_0 = 1$。$m$ 称为该递推式的 \markdownRendererStrongEmphasis{阶数}。\markdownRendererParagraphSeparator
{}数列 ${a_i}$ 的最短递推式即为阶数最小的递推式。\markdownRendererInterblockSeparator
{}
\markdownRendererSectionEnd \markdownRendererSectionBegin
\markdownRendererHeadingThree{做法}\markdownRendererInterblockSeparator
{}与上面定义的稍有不同，这里定义一个新的递推系数 ${f_0 \dots f_{m - 1}}$，满足：\markdownRendererParagraphSeparator
{}$a_i = \sum_{j = 0} ^ {m - 1} f_j a_{i - j - 1}, \forall i \ge m$\markdownRendererParagraphSeparator
{}容易看出 $f_i = -r_{i + 1}$，并且阶数 $m$ 与之前的定义是相同的。\markdownRendererParagraphSeparator
{}我们可以增量地求递推式，按顺序考虑 ${a_i}$ 的每一位，并在递推结果出现错误时对递推系数 ${f_i}$ 进行调整。方便起见，以下将前 $i$ 位的最短递推式记为 $F_i = {f_{i, j}}$。\markdownRendererParagraphSeparator
{}显然初始时有 $F_0 = {}$。假设递推系数 $F_{i - 1}$ 对数列 ${a_i}$ 的前 $i - 1$ 项均成立，这时对第 $i$ 项就有两种情况：\markdownRendererInterblockSeparator
{}\markdownRendererOlBeginTight
\markdownRendererOlItemWithNumber{1}递推系数对 $a_i$ 也成立，这时不需要进行任何调整，直接令 $F_i = F_{i - 1}$ 即可。\markdownRendererOlItemEnd 
\markdownRendererOlItemWithNumber{2}递推系数对 $a_i$ 不成立，这时需要对 $F_{i - 1}$ 进行调整，得到新的 $F_i$。\markdownRendererOlItemEnd 
\markdownRendererOlEndTight \markdownRendererInterblockSeparator
{}设 $\Delta_i = a_i - \sum_{j = 0} ^ m f_{i - 1, j} a_{i - j - 1}$，即 $a_i$ 与 $F_{i - 1}$ 的递推结果的差值。\markdownRendererParagraphSeparator
{}如果这是第一次对递推系数进行修改，则说明 $a_i$ 是序列中的第一个非零项。这时直接令 $F_i$ 为 $i$ 个 $0$ 即可，显然这是一个合法的最短递推式。\markdownRendererParagraphSeparator
{}否则设上一次对递推系数进行修改时，已考虑的 ${a_i}$ 的项数为 $k$。如果存在一个序列 $G = {g_0 \dots g_{m' - 1}}$，满足：\markdownRendererParagraphSeparator
{}$\sum_{j = 0} ^ {m' - 1} g_j a_{i' - j - 1} = 0, \forall i' \in [m', i)$\markdownRendererParagraphSeparator
{}并且 $\sum_{j = 0} ^ {m' - 1} g_j a_{i - j - 1} = \Delta_i$，那么不难发现将 $F_k$ 与 $G$ 按位分别相加之后即可得到一个合法的递推系数 $F_i$。\markdownRendererParagraphSeparator
{}考虑如何构造 $G$。一种可行的构造方案是令\markdownRendererParagraphSeparator
{}$G = {0, 0, \dots, 0, \frac{\Delta_i}{\Delta_k}, -\frac{\Delta_i}{\Delta_k}F_{k-1}}$\markdownRendererParagraphSeparator
{}其中前面一共有 $i - k - 1$ 个 $0$，且最后的 $-\frac{\Delta_i}{\Delta_k} F_{k-1}$ 表示将 $F_{k-1}$ 每项乘以 $-\frac{\Delta_i}{\Delta_k}$ 后接在序列后面。\markdownRendererParagraphSeparator
{}不难验证此时 $\sum_{j = 0} ^ {m' - 1} g_j a_{i - j - 1} = \Delta_k \frac{\Delta_i}{\Delta_k} = \Delta_i$，因此这样构造出的是一个合法的 $G$。将 $F_i$ 赋值为 $F_k$ 与 $G$ 逐项相加后的结果即可。\markdownRendererParagraphSeparator
{}如果要求的是符合最开始定义的递推式 ${r_i}$，则将 ${f_j}$ 全部取相反数后在最开始插入 $r_0 = 1$ 即可。\markdownRendererParagraphSeparator
{}从上述算法流程中可以看出，如果数列的最短递推式的阶数为 $m$，则算法的复杂度为 $O(nm)$。最坏情况下 $m = O(n)$，因此算法的最坏复杂度为 $O(n^2)$。\markdownRendererParagraphSeparator
{}在实现算法时，由于每次调整递推系数时都只需要用到上次调整时的递推系数 $F_k$，因此如果只需要求整个数列的最短递推式，可以只存储当前递推系数和上次调整时的递推系数，空间复杂度为 $O(n)$。\markdownRendererInterblockSeparator
{}
\markdownRendererSectionEnd \markdownRendererSectionBegin
\markdownRendererHeadingThree{参考实现}\markdownRendererInterblockSeparator
{}\markdownRendererInputFencedCode{./_markdown_main/147468f5cb267ec2ce60e67a48ef5d9b.verbatim}{cpp}{cpp}\markdownRendererInterblockSeparator
{}朴素的 Berlekamp–Massey 算法求解的是有限项数列的最短递推式。如果待求递推式的序列有无限项，但已知最短递推式的阶数上界，则只需取出序列的前 $2m$ 项即可求出整个序列的最短递推式。（证明略）\markdownRendererInterblockSeparator
{}
\markdownRendererSectionEnd \markdownRendererSectionBegin
\markdownRendererHeadingThree{应用}\markdownRendererInterblockSeparator
{}由于 Berlekamp–Massey 算法的数值稳定性比较差，在处理实数问题时一般很少使用。为了叙述方便，以下均假定在某个质数 $p$ 的剩余系下进行运算。\markdownRendererInterblockSeparator
{}\markdownRendererSectionBegin
\markdownRendererHeadingFour{求向量列或矩阵列的最短递推式}\markdownRendererInterblockSeparator
{}如果要求向量列 $\boldsymbol{v}_i$ 的最短递推式，设向量的维数为 $n$，我们可以随机一个 $n$ 维行向量 $\mathbf u^T$，并计算标量序列 ${\boldsymbol{u}^T\boldsymbol{v}_i}$ 的最短递推式。由 Schwartz–Zippel 引理，二者的最短递推式有至少 $1 - \frac n p$ 的概率相同。\markdownRendererParagraphSeparator
{}求矩阵列 ${A_i}$ 的最短递推式也是类似的，设矩阵的大小为 $n \times m$，则只需随机一个 $1 \times n$ 的行向量 $\mathbf u^T$ 和一个 $m \times 1$ 的列向量 $\boldsymbol{v}$，并计算标量序列 ${\boldsymbol{u}^T A_i \boldsymbol{v}}$ 的最短递推式即可。由 Schwartz–Zippel 引理可以类似地得到二者相同的概率至少为 $1 - \frac{n + m} p$。\markdownRendererInterblockSeparator
{}
\markdownRendererSectionEnd \markdownRendererSectionBegin
\markdownRendererHeadingFour{优化矩阵快速幂}\markdownRendererInterblockSeparator
{}设 $\boldsymbol{f}_i$ 是一个 $n$ 维列向量，并且转移满足 $\boldsymbol{f}\markdownRendererEmphasis{i = A \boldsymbol{f}}{i-1}$（其中下标 $i-1$），则可以发现 ${\boldsymbol{f}_i}$ 是一个不超过 $n$ 阶的线性递推向量列。（证明略）\markdownRendererParagraphSeparator
{}我们可以直接暴力求出 $\boldsymbol{f}\markdownRendererEmphasis{0 \dots \boldsymbol{f}}{2n - 1}$，然后用前面提到的做法求出 ${\boldsymbol{f}_i}$ 的最短递推式，再调用 \markdownRendererLink{常系数齐次线性递推}{./poly/linear-recurrence.md}{./poly/linear-recurrence.md}{} 即可。\markdownRendererParagraphSeparator
{}如果要求的向量是 $\boldsymbol{f}_m$，则算法的复杂度是 $O(n^3 + n\log n \log m)$。如果 $A$ 是一个只有 $k$ 个非零项的稀疏矩阵，则复杂度可以降为 $O(nk + n\log n \log m)$。但由于算法至少需要 $O(nk)$ 的时间预处理，因此在压力不大的情况下也可以使用 $O(n^2 \log m)$ 的线性递推算法，复杂度同样是可以接受的。\markdownRendererInterblockSeparator
{}
\markdownRendererSectionEnd \markdownRendererSectionBegin
\markdownRendererHeadingFour{求矩阵的最小多项式}\markdownRendererInterblockSeparator
{}方阵 $A$ 的最小多项式是次数最小的并且满足 $f(A) = 0$ 的多项式 $f$。\markdownRendererParagraphSeparator
{}实际上最小多项式就是 ${A^i}$ 的最小递推式，所以直接调用 Berlekamp–Massey 算法就可以了。如果 $A$ 是一个 $n$ 阶方阵，则显然最小多项式的次数不超过 $n$。\markdownRendererParagraphSeparator
{}瓶颈在于求出 $A^i$，因为如果直接每次做矩阵乘法的话复杂度会达到 $O(n^4)$。但考虑到求矩阵列的最短递推式时实际上求的是 ${\boldsymbol{u}^T A^i \boldsymbol{v}}$ 的最短递推式，因此我们只要求出 $A^i \boldsymbol{v}$ 就行了。\markdownRendererParagraphSeparator
{}假设 $A$ 有 $k$ 个非零项，则复杂度为 $O(kn + n^2)$。\markdownRendererInterblockSeparator
{}
\markdownRendererSectionEnd \markdownRendererSectionBegin
\markdownRendererHeadingFour{求稀疏矩阵行列式}\markdownRendererInterblockSeparator
{}如果能求出方阵 $A$ 的特征多项式，则常数项乘上 $(-1)^n$ 就是行列式。但是最小多项式不一定就是特征多项式。\markdownRendererParagraphSeparator
{}实际上如果把 $A$ 乘上一个随机对角阵 $B$，则 $AB$ 的最小多项式有至少 $1 - \frac {2n^2 - n} p$ 的概率就是特征多项式。最后再除掉 $\text{det};B$ 就行了。\markdownRendererParagraphSeparator
{}设 $A$ 为 $n$ 阶方阵，且有 $k$ 个非零项，则复杂度为 $O(kn + n ^ 2)$。\markdownRendererInterblockSeparator
{}
\markdownRendererSectionEnd \markdownRendererSectionBegin
\markdownRendererHeadingFour{求稀疏矩阵的秩}\markdownRendererInterblockSeparator
{}设 $A$ 是一个 $n\times m$ 的矩阵，首先随机一个 $n\times n$ 的对角阵 $P$ 和一个 $m\times m$ 的对角阵 $Q$, 然后计算 $Q A P A^T Q$ 的最小多项式即可。\markdownRendererParagraphSeparator
{}实际上不用调用矩阵乘法，因为求最小多项式时要用 $Q A P A^T Q$ 乘一个向量，所以我们依次把这几个矩阵乘到向量里就行了。答案就是最小多项式除掉所有 $x$ 因子后剩下的次数。\markdownRendererParagraphSeparator
{}设 $A$ 有 $k$ 个非零项，且 $n \le m$，则复杂度为 $O(kn + n ^ 2)$。\markdownRendererInterblockSeparator
{}
\markdownRendererSectionEnd \markdownRendererSectionBegin
\markdownRendererHeadingFour{解稀疏方程组}\markdownRendererInterblockSeparator
{}\markdownRendererStrongEmphasis{问题}：已知 $A \mathbf x = \mathbf b$, 其中 $A$ 是一个 $n \times n$ 的 \markdownRendererStrongEmphasis{满秩} 稀疏矩阵，$\mathbf b$ 和 $\mathbf x$ 是 $1\times n$ 的列向量。$A, \mathbf b$ 已知，需要在低于 $n^\omega$ 的复杂度内解出 $x$。\markdownRendererParagraphSeparator
{}\markdownRendererStrongEmphasis{做法}：显然 $\mathbf x = A^{-1} \mathbf b$。如果我们能求出 ${A^i \mathbf b}$($i \ge 0$) 的最小递推式 ${r_0 \dots r_{m - 1}}$($m \le n$), 那么就有结论\markdownRendererParagraphSeparator
{}$A^{-1} \mathbf b = -\frac 1 {r_{m - 1}} \sum_{i = 0} ^ {m - 2} A^i \mathbf b r_{m - 2 - i}$\markdownRendererParagraphSeparator
{}（证明略）\markdownRendererParagraphSeparator
{}因为 $A$ 是稀疏矩阵，直接按定义递推出 $\mathbf b \dots A^{2n - 1} \mathbf b$ 即可。\markdownRendererParagraphSeparator
{}同样地，设 $A$ 中有 $k$ 个非零项，则复杂度为 $O(kn + n^2)$。\markdownRendererInterblockSeparator
{}
\markdownRendererSectionEnd 
\markdownRendererSectionEnd \markdownRendererSectionBegin
\markdownRendererHeadingThree{参考实现}\markdownRendererInterblockSeparator
{}\markdownRendererInputFencedCode{./_markdown_main/f177f28619689a0e25eadc243523dfae.verbatim}{cpp}{cpp}\markdownRendererInterblockSeparator
{}
\markdownRendererSectionEnd \markdownRendererSectionBegin
\markdownRendererHeadingThree{例题}\markdownRendererInterblockSeparator
{}\markdownRendererOlBeginTight
\markdownRendererOlItemWithNumber{1}\markdownRendererLink{LibreOJ #163. 高斯消元 2}{https://loj.ac/p/163}{https://loj.ac/p/163}{}\markdownRendererOlItemEnd 
\markdownRendererOlItemWithNumber{2}\markdownRendererLink{ICPC2021 台北 Gym103443E. Composition with Large Red Plane, Yellow, Black, Gray, and Blue}{https://codeforces.com/gym/103443/problem/E}{https://codeforces.com/gym/103443/problem/E}{}\markdownRendererOlItemEnd 
\markdownRendererOlEndTight 
\markdownRendererSectionEnd 
\markdownRendererSectionEnd 
\markdownRendererSectionEnd \markdownRendererDocumentEnd